% Generado a partir de data/huellas/vivo.json y dashboard/data/inundacion.json (recálculo del 4/11/2026)
\newcommand{\VivoEstado}{\hecho}
\newcommand{\VivoTexto}{%
La capa de agua nueva de los últimos 24 días se recalculó con Sentinel-1 después de corregir el error E4.
Usó una referencia seca del 1/8 al 10/10/2026 (42 pasadas) y un evento del 11/10 al 4/11/2026 (13 pasadas, del 12/10 al 2/11).
\begin{itemize}
\item \textbf{Detección:} 334,8 km² en bruto dentro de la caja de estudio; quedan \textbf{27,2 km²} en tierra de Tumbes y Piura después de aplicar las máscaras. El 77\,\% está en Sechura.
\item \textbf{Exposición:} 24 personas y 8 edificaciones dentro de la mancha; ningún colegio ni establecimiento de salud dentro.
\item \textbf{Umbral:} el umbral automático de Otsu, calculado como diagnóstico, fue $-16{,}8$~dB, muy cercano al umbral fijo de $-17$~dB.
\end{itemize}
Es temporada seca, así que el agua nueva corresponde sobre todo a riego y humedales, no a desbordes. La capa pasa el control de calidad y su estado es \emph{vigente}.}
